\section{Алгебра событий}
\subsection{Основные определения}
\subsubsection{События. Исходы. Пространство элементарных исходов}

Для начала следует познакомиться с тремя основополагающими понятиями в Алгебре событий:

\begin{definition}[Элементарный исход]
\textbf{Элементарный исход} (элементарное событие, элементарный исход эксперимента) -- это отдельный, далее неделимый результат случайного эксперимента, обозначаемый как $\omega \in \Omega$. Каждый элементарный исход представляет собой один конкретный вариант того, чем может закончиться эксперимент.
\end{definition}

\begin{definition}[Событие]
\textbf{Событие} -- множество, состоящее из элементарных исходов данного эксперимента, т.е. некоторое подмножество пространства элементарных исходов: $A \subseteq \Omega$. Событие происходит (наступает), если в результате эксперимента реализуется хотя бы один элементарный исход $\omega$, принадлежащий этому множеству.
\end{definition}

\begin{definition}[Пространство элементарных исходов]
Пространство элементарных исходов/событий -- множество $\Omega$ событий (всех возможных исходов), являющихся исходами данного случайного эксперимента, из которых в эксперименте происходит ровно один.
\end{definition}

\textbf{Пример:}
    \center\textit{Эксперимент:} бросание игральной кости.
    \begin{itemize}
    \item Пространство элементарных исходов: $\Omega = \{1, 2, 3, 4, 5, 6\}$.
    
    \item Элементарный исход: $\omega = 4$ (выпала четвёрка).
    
    \item Событие $A$: ``выпало число больше 4'' $\Rightarrow A = \{5, 6\}$.
    
    \item Событие $B$: ``выпало нечётное число'' $\Rightarrow B = \{1, 3, 5\}$.
    \end{itemize}

\begin{notice}Таким образом:

\begin{enumerate}
    \item Элементарный исход -- это \textit{элемент} множества $\Omega$ (обозначается $\omega \in \Omega$).
    \item Событие -- это \textit{подмножество} $\Omega$ (обозначается $A \subseteq \Omega$).
    \item Событие может состоять из одного исхода (такое событие называется \textit{элементарным}) или из нескольких исходов.
    \item Всё пространство $\Omega$ тоже является событием (достоверное событие), равно как и пустое множество $\varnothing$ -- тоже событие (невозможное событие).
\end{enumerate}
\end{notice}

Элементарные исходы, которые являеются элементами данного события,
называют \textbf{благоприятными}.

\subsubsection{Операции над событиями}

Поскольку события мы определили как множества, то событиям свойственны множественные операции

\begin{definition}[Пересечение (произведение) событий]
Событие $C$, происходящее тогда и только тогда, когда наступает и событие $A$ и событие $B$, называется \textbf{пересечением (произведением) событий $A$ и $B$}. Обозначение: $A \cap B$ или $AB$.
\end{definition}

\begin{definition}[Несовместные события]
События $A$ и $B$ называют \textbf{несовместными}, если они не могут произойти одновременно, т.е. $A \cap B = \varnothing$.
\end{definition}

В противном случае события называют совместными или пересекающимися.

\begin{definition}[Объединение (сумма) событий]
Событие $C$, происходящее тогда и только тогда, когда происходит хотя бы одно из событий $A$ или $B$, называется \textbf{объединением (суммой) событий $A$ и $B$}. Обозначение: $A \cup B$ или $A + B$.
\end{definition}